\begin{enumerate}

\item[\textbf{آ)}]

در این سؤال ضریب تخفیف برابر است با:

\[
\gamma = 0.5
\]

سیاست اولیه \(\pi_1\) در هر حالت، بین دو کنش \(x\) و \(y\) با احتمال برابر انتخاب می‌کند. همچنین مقدار اولیه‌ی تابع ارزش برای همه‌ی حالت‌ها برابر است با:

\[
V_1(A)=V_1(B)=V_1(C)=2
\]

برای اجرای یک مرحله از ارزیابی سیاست، مقدار هر حالت را با استفاده از مقدارهای قبلی حساب می‌کنیم.

\medskip

\textbf{حالت \(A\):}

اگر کنش \(x\) را انتخاب کنیم، با پاداش \(4\) به حالت \(B\) می‌رویم. پس:

\[
Q_1(A,x)=4+\gamma V_1(B)=4+0.5(2)=5
\]

اگر کنش \(y\) را انتخاب کنیم، با پاداش \(2\) به حالت \(C\) می‌رویم. پس:

\[
Q_1(A,y)=2+\gamma V_1(C)=2+0.5(2)=3
\]

چون سیاست تصادفی است، مقدار جدید \(A\) برابر میانگین این دو مقدار می‌شود:

\[
V_2(A)=\frac{1}{2}Q_1(A,x)+\frac{1}{2}Q_1(A,y)
=
\frac{1}{2}(5)+\frac{1}{2}(3)=4
\]

\medskip

\textbf{حالت \(B\):}

اگر کنش \(x\) را انتخاب کنیم، با پاداش \(2\) به حالت \(C\) می‌رویم:

\[
Q_1(B,x)=2+\gamma V_1(C)=2+0.5(2)=3
\]

اگر کنش \(y\) را انتخاب کنیم، با پاداش \(0\) به حالت \(A\) می‌رویم:

\[
Q_1(B,y)=0+\gamma V_1(A)=0+0.5(2)=1
\]

پس:

\[
V_2(B)=\frac{1}{2}Q_1(B,x)+\frac{1}{2}Q_1(B,y)
=
\frac{1}{2}(3)+\frac{1}{2}(1)=2
\]

\medskip

\textbf{حالت \(C\):}

اگر کنش \(y\) را انتخاب کنیم، با پاداش \(1\) به حالت \(B\) می‌رویم:

\[
Q_1(C,y)=1+\gamma V_1(B)=1+0.5(2)=2
\]

اگر کنش \(x\) را انتخاب کنیم، پاداش \(8\) می‌گیریم و سپس با احتمال \(\frac{1}{4}\) به \(A\) می‌رویم و با احتمال \(\frac{3}{4}\) در \(C\) می‌مانیم. بنابراین:

\[
Q_1(C,x)
=
8+\gamma\left(\frac{1}{4}V_1(A)+\frac{3}{4}V_1(C)\right)
\]

\[
Q_1(C,x)
=
8+0.5\left(\frac{1}{4}(2)+\frac{3}{4}(2)\right)
=
8+0.5(2)=9
\]

پس:

\[
V_2(C)=\frac{1}{2}Q_1(C,x)+\frac{1}{2}Q_1(C,y)
=
\frac{1}{2}(9)+\frac{1}{2}(2)=5.5
\]

در نتیجه بعد از یک مرحله ارزیابی سیاست داریم:

\[
\boxed{
V_2(A)=4,\qquad V_2(B)=2,\qquad V_2(C)=5.5
}
\]

\item[\textbf{ب)}]

حالا باید نسبت به تابع ارزش جدید، یعنی \(V_2\)، سیاست حریصانه‌ی جدید \(\pi_2\) را پیدا کنیم.

\[
V_2(A)=4,\qquad V_2(B)=2,\qquad V_2(C)=5.5
\]

\medskip

\textbf{حالت \(A\):}

\[
Q_2(A,x)=4+\gamma V_2(B)=4+0.5(2)=5
\]

\[
Q_2(A,y)=2+\gamma V_2(C)=2+0.5(5.5)=4.75
\]

چون:

\[
Q_2(A,x)>Q_2(A,y)
\]

پس:

\[
\pi_2(A)=x
\]

\medskip

\textbf{حالت \(B\):}

\[
Q_2(B,x)=2+\gamma V_2(C)=2+0.5(5.5)=4.75
\]

\[
Q_2(B,y)=0+\gamma V_2(A)=0+0.5(4)=2
\]

چون:

\[
Q_2(B,x)>Q_2(B,y)
\]

پس:

\[
\pi_2(B)=x
\]

\medskip

\textbf{حالت \(C\):}

\[
Q_2(C,y)=1+\gamma V_2(B)=1+0.5(2)=2
\]

\[
Q_2(C,x)
=
8+\gamma\left(\frac{1}{4}V_2(A)+\frac{3}{4}V_2(C)\right)
\]

\[
Q_2(C,x)
=
8+0.5\left(\frac{1}{4}(4)+\frac{3}{4}(5.5)\right)
\]

\[
Q_2(C,x)
=
8+0.5(1+4.125)
=
8+2.5625
=
10.5625
\]

چون:

\[
Q_2(C,x)>Q_2(C,y)
\]

پس:

\[
\pi_2(C)=x
\]

در نتیجه سیاست حریصانه‌ی جدید برابر است با:

\[
\boxed{
\pi_2(A)=x,\qquad \pi_2(B)=x,\qquad \pi_2(C)=x
}
\]

\item[\textbf{ج)}]

می‌خواهیم نشان دهیم اگر سیاست \(\pi'\) به صورت حریصانه از روی \(V^\pi\) ساخته شده باشد، آنگاه برای همه‌ی حالت‌ها داریم:

\[
V^{\pi'}(s)\ge V^\pi(s)
\]

چون \(\pi'\) نسبت به \(V^\pi\) حریصانه است، برای هر حالت \(s\) داریم:

\[
\pi'(s)=\arg\max_a Q^\pi(s,a)
\]

پس:

\[
Q^\pi(s,\pi'(s))=\max_a Q^\pi(s,a)
\]

از طرفی کنشی که سیاست \(\pi\) انتخاب می‌کند، یکی از همین کنش‌هاست. بنابراین:

\[
Q^\pi(s,\pi'(s))\ge Q^\pi(s,\pi(s))
\]

ولی طبق تعریف تابع ارزش سیاست \(\pi\) داریم:

\[
Q^\pi(s,\pi(s))=V^\pi(s)
\]

پس نتیجه می‌شود:

\[
Q^\pi(s,\pi'(s))\ge V^\pi(s)
\]

حالا \(Q^\pi(s,\pi'(s))\) را با رابطه‌ی بلمن باز می‌کنیم:

\[
Q^\pi(s,\pi'(s))
=
R(s,\pi'(s))
+
\gamma
\sum_{s'}P(s'|s,\pi'(s))V^\pi(s')
\]

پس داریم:

\[
V^\pi(s)
\le
R(s,\pi'(s))
+
\gamma
\sum_{s'}P(s'|s,\pi'(s))V^\pi(s')
\]

حالا اگر همین نامساوی را برای حالت‌های بعدی هم به صورت تکراری جایگذاری کنیم، سمت راست به ارزش اجرای سیاست \(\pi'\) تبدیل می‌شود. یعنی:

\[
V^\pi(s)
\le
\mathbb{E}_{\pi'}\left[
\sum_{t=0}^{\infty}\gamma^t R_t
\mid S_0=s
\right]
\]

پس:

\[
V^\pi(s)\le V^{\pi'}(s)
\]

بنابراین:

\[
\boxed{
V^{\pi'}(s)\ge V^\pi(s)
}
\]

برای بخش دوم، فرض کنیم برای همه‌ی حالت‌ها تساوی برقرار باشد:

\[
V^{\pi'}(s)=V^\pi(s)
\]

از طرفی چون \(\pi'\) نسبت به \(V^\pi\) حریصانه ساخته شده است، داریم:

\[
T^*V^\pi = T^{\pi'}V^\pi
\]

و چون طبق فرض \(V^{\pi'}=V^\pi\)، پس:

\[
V^\pi = T^{\pi'}V^\pi
\]

بنابراین:

\[
V^\pi = T^*V^\pi
\]

این یعنی \(V^\pi\) معادله‌ی بهینگی بلمن را ارضا می‌کند. از آن‌جا که جواب معادله‌ی بهینگی بلمن یکتاست، داریم:

\[
V^\pi = V^*
\]

پس سیاست \(\pi\) بهینه است. همچنین چون \(\pi'\) نسبت به \(V^*\) حریصانه است، \(\pi'\) هم سیاست بهینه خواهد بود.

در نتیجه اگر تساوی برای همه‌ی حالت‌ها رخ دهد، سیاست به دست آمده حتماً بهینه است.

\[
\boxed{
V^{\pi'}=V^\pi \Longrightarrow V^\pi=V^*
}
\]

\end{enumerate}